\subsection{MLOps y l\'ogica de umbrales}
Se propuso una pol\'itica operacional basada en la probabilidad posterior m\'axima del modelo Soft Voting Ensemble. Las predicciones con $P \geq 0.85$ fueron asignadas a la zona de confianza, las predicciones con $0.40 \leq P < 0.85$ fueron asignadas a la zona de incertidumbre y las predicciones con $P < 0.40$ fueron asignadas a la zona de rechazo.

\begin{table}[H]
\centering
\caption{Resumen operativo de zonas de decisi\'on.}
\begin{tabular}{lrrrr}
\hline
Zona & Observaciones & Cobertura & Accuracy & Prob. media \\
\hline
Zona de Confianza & 9 & 0.019 & 0.667 & 0.883 \\
Zona de Incertidumbre & 417 & 0.874 & 0.492 & 0.588 \\
Zona de Rechazo & 51 & 0.107 & 0.451 & 0.353 \\
\hline
\end{tabular}
\end{table}

La zona de confianza fue definida para automatizar decisiones de clasificaci\'on con alta certeza. La zona de incertidumbre fue destinada a revisi\'on humana o validaci\'on secundaria, mientras que la zona de rechazo fue reservada para casos donde no deber\'ia emitirse una clasificaci\'on autom\'atica. Esta l\'ogica permite controlar el riesgo operativo al separar cobertura del sistema y confiabilidad de las decisiones aceptadas.
